\chapter{论文地图 Paper Map}
本附录按章节重排 supplemental readings，帮助读者把 lecture 主线与官方 papers / blog / project page 对齐。
\small
\begin{longtable}{p{0.11\textwidth}p{0.25\textwidth}p{0.24\textwidth}p{0.3\textwidth}}
\toprule
章节 & 论文 / Reading & 主要问题 & 与本章的连接 \\
\midrule
\endhead
Chapter 1 & Large Language Models as Optimizers\newline\url{https://arxiv.org/abs/2309.03409} & Can an LLM act as an optimizer by proposing candidate instructions or solutions from the history of past trials and scores? & This is the lecture's canonical example for turning prompt design into an explicit inference-time search loop over instructions. \\
Chapter 1 & Large Language Models Cannot Self-Correct Reasoning Yet\newline\url{https://arxiv.org/abs/2310.01798} & Can an LLM improve reasoning quality by revising its own answer without any reliable external feedback? & This reading directly grounds the lecture's warning that iterative self-improvement only works when the feedback signal is good. \\
Chapter 1 & Teaching Large Language Models to Self-Debug\newline\url{https://arxiv.org/abs/2304.05128} & Can LLMs debug their own generated programs more effectively when they are asked to explain, test, and revise code? & This is the lecture's strongest positive case for iterative self-improvement, because code tasks expose external signals such as tests and traces. \\
Chapter 2 & Direct Preference Optimization: Your Language Model is Secretly a Reward Model\newline\url{https://arxiv.org/abs/2305.18290} & Can we replace RLHF's separate reward-model-plus-RL stack with a simpler objective that directly optimizes pairwise preferences? & Jason Weston uses DPO as the main optimization primitive inside self-rewarding and iterative reasoning pipelines, so this paper is the lecture's optimizer backbone. \\
Chapter 2 & Iterative Reasoning Preference Optimization\newline\url{https://arxiv.org/abs/2404.19733} & How can preference optimization be adapted so that it actually improves difficult reasoning tasks rather than only generic instruction following? & This paper is the lecture's central answer to the question 'how do we teach reasoning rather than only obedience or style alignment?' \\
Chapter 2 & Chain-of-Verification Reduces Hallucination in Large Language Models\newline\url{https://arxiv.org/abs/2309.11495} & Can a language model reduce hallucination by explicitly planning verification questions and answering them independently before finalizing its response? & The lecture uses CoVe to motivate why System 2 style reasoning should often be framed as explicit verification rather than just longer free-form chain-of-thought. \\
Chapter 3 & Grokked Transformers are Implicit Reasoners: A Mechanistic Journey to the Edge of Generalization\newline\url{https://arxiv.org/abs/2405.15071} & Can transformers learn implicit reasoning over parametric knowledge, and what happens internally when they eventually generalize? & Yu Su uses this paper to argue that language-agent reasoning is not only about explicit chain-of-thought; implicit parametric reasoning remains a core substrate. \\
Chapter 3 & HippoRAG: Neurobiologically Inspired Long-Term Memory for Large Language Models\newline\url{https://arxiv.org/abs/2405.14831} & How can LLM memory retrieval move beyond shallow semantic similarity and support deeper knowledge integration over long-term experiences? & This is Yu Su's flagship example for why language agents need memory architectures richer than vanilla retrieval-augmented generation. \\
Chapter 3 & Is Your LLM Secretly a World Model of the Internet? Model-Based Planning for Web Agents\newline\url{https://arxiv.org/abs/2411.06559} & Can web agents plan by simulating future website states with a world model instead of relying only on reactive execution or expensive tree search? & Yu Su uses WebDreamer to show how language-agent planning becomes tractable when LLMs can predict environment transitions before acting. \\
Chapter 4 & Tulu 3: Pushing Frontiers in Open Language Model Post-Training\newline\url{https://arxiv.org/abs/2411.15124} & What fully-open recipe can push post-training performance of open models close to or beyond strong proprietary and open instruct baselines? & This is the canonical paper behind the lecture's three-stage recipe: SFT, preference tuning, and RLVR. \\
Chapter 4 & Unpacking DPO and PPO: Disentangling Best Practices for Learning from Preference Feedback\newline\url{https://arxiv.org/abs/2406.09279} & When doing preference-based post-training, which ingredients matter most: data, algorithm, reward model, or prompt construction? & This reading grounds Hanna Hajishirzi's message that algorithm choice matters, but data and evaluation quality matter even more. \\
Chapter 4 & OpenScholar: Synthesizing Scientific Literature with Retrieval-augmented LMs\newline\url{https://arxiv.org/abs/2411.14199} & Can an open retrieval-augmented language model produce citation-grounded scientific synthesis at a quality level competitive with stronger closed systems? & The lecture is mostly about training recipes, but OpenScholar is an important downstream example of why open infrastructure and reproducible open models matter. \\
Chapter 5 & EnIGMA: Interactive Tools Substantially Assist LM Agents in Finding Security Vulnerabilities\newline\url{https://arxiv.org/abs/2409.16165} & How much do interactive tools such as debuggers and server-connection utilities improve security-oriented LM agents on CTF tasks? & This is the lecture's clearest example of why security agents need richer tools than ordinary coding agents. \\
Chapter 5 & From Naptime to Big Sleep: Using Large Language Models To Catch Vulnerabilities In Real-World Code\newline\url{https://projectzero.google/2024/10/from-naptime-to-big-sleep.html} & Can an LLM agent find previously unknown exploitable vulnerabilities in real-world code before release? & This blog post anchors the lecture's final case study on real-world vulnerability detection and explains why tooling plus verification matter more than generic 'chat intelligence'. \\
Chapter 6 & Mind2Web: Towards a Generalist Agent for the Web\newline\url{https://arxiv.org/abs/2306.06070} & 如何构造一个能够在真实网站上执行开放式任务、并能泛化到未见网站与未见 domain 的 generalist web agent 数据集与基线？ & 本讲把它当作 web agent 数据与 benchmark 家族中的“真实网站、离线数据集”基线，对比后续 VisualWebArena 的视觉 grounding 与 Tree Search 的 inference-time enhancement。 \\
Chapter 6 & WebArena: A Realistic Web Environment for Building Autonomous Agents\newline\url{https://arxiv.org/abs/2307.13854} & 怎样构造一个 realistic yet reproducible 的 web environment，用来评估 language-guided autonomous agents 在真实互联网任务上的 functional correctness？ & Ruslan 把它作为 VisualWebArena 的前身和对照组，用来说明“仅靠 HTML 不够，真实网页操作需要视觉输入”。 \\
Chapter 6 & VisualWebArena: Evaluating Multimodal Agents on Realistic Visual Web Tasks\newline\url{https://jykoh.com/vwa} & 如果 web tasks 需要图像、布局和视觉线索，如何系统评估 multimodal web agents 的真实表现？ & 这是本讲前半段的主线 benchmark，支撑“视觉 grounding 对 web agents 是必要条件，而非锦上添花”。 \\
Chapter 6 & Tree Search for Language Model Agents\newline\url{https://jykoh.com/search-agents} & 面对 realistic web tasks 中的长程规划和局部决策误差累积，能否通过 inference-time search 改善 language model agent？ & 这是本讲中“multimodal agent 不是只能靠更强 base model，还能靠 inference-time compute”最直接的例子。 \\
Chapter 7 & OSWORLD: Benchmarking Multimodal Agents for Open-Ended Tasks in Real Computer Environments\newline\url{https://arxiv.org/pdf/2404.07972} & 如何为 multimodal computer-use agents 构造第一个可扩展、真实、可执行评测的计算机环境和 benchmark？ & 这是本讲前半段的 canonical benchmark reading，用来定义真实 computer-use agent 的评测对象。 \\
Chapter 7 & AGUVIS: Unified Pure Vision Agents for Autonomous GUI Interaction\newline\url{https://arxiv.org/pdf/2412.04454} & 能否构造一种 purely vision-based、跨平台统一的 GUI agent，不再依赖 HTML/AXTree/XML 等异构文本表示？ & 这是本讲中“从 perception 到 action”的核心 reading，直接支撑 GUI grounding、cross-platform action space 和 inner monologue sections。 \\
Chapter 8 & Mastering Chess and Shogi by Self-Play with a General Reinforcement Learning Algorithm\newline\url{https://arxiv.org/abs/1712.01815} & Can a single self-play reinforcement learning recipe discover superhuman strategies across games without handcrafted evaluation functions? & L08 reinterprets Lean as the mathematical analogue of a game environment and treats AlphaProof as an AlphaZero-style agent for proof search rather than board play. \\
Chapter 8 & The Future of Mathematics?\newline\url{https://www.youtube.com/watch?v=Dp-mQ3HxgDE} & How could theorem provers such as Lean change mathematical practice, education, and eventually research? & Thomas Hubert's lecture inherits this ecosystem view: AlphaProof only makes sense because Lean, Mathlib, and the formalization community are turning mathematics into a machine-verifiable environment. \\
Chapter 9 & LeanDojo: Theorem Proving with Retrieval-Augmented Language Models\newline\url{https://arxiv.org/abs/2306.15626} & How can researchers build an open, reproducible training and evaluation environment for Lean theorem proving with retrieval-augmented language models? & Kaiyu Yang uses LeanDojo as the canonical example of theorem proving as an interaction problem between a language model, a proof environment, and a retrieval system. \\
Chapter 9 & Autoformalization with Large Language Models\newline\url{https://arxiv.org/abs/2205.12615} & Can large language models translate natural-language mathematics into formal specifications and proofs well enough to improve downstream theorem proving? & The lecture uses this paper to define autoformalization as a distinct task from theorem proving: the system must first write the right formal statement before it can hope to prove it. \\
Chapter 9 & Autoformalizing Euclidean Geometry\newline\url{https://arxiv.org/abs/2405.17216} & How can we autoformalize Euclidean geometry when informal proofs depend heavily on diagrams and unstated geometric constraints? & This paper grounds the lecture's argument that autoformalization is hardest precisely where informal proofs rely on hidden diagrammatic reasoning. \\
Chapter 10 & Draft, Sketch, and Prove: Guiding Formal Theorem Provers with Informal Proofs\newline\url{https://arxiv.org/abs/2210.12283} & 如何把 informal proof 变成 formal prover 的有效搜索偏置，而不是把自然语言解释直接当成最终证明。 & 这是 lecture 中“combining informal and formal provers”的核心 reading，直接支撑第 3 节。 \\
Chapter 10 & miniCTX: Neural Theorem Proving with (Long-)Contexts\newline\url{https://arxiv.org/abs/2408.03350} & 如何评估 theorem prover 是否能利用真实项目中的长上下文，而不是只会解短小自包含题。 & 它对应本讲最后一部分对 research-level formalization 的转向，强调真实项目环境与竞赛题环境的差异。 \\
Chapter 10 & Lean-STaR: Learning to Interleave Thinking and Proving\newline\url{https://arxiv.org/abs/2407.10040} & formal proof data 并不显示人类的思考过程，能否在 tactic 前显式学习 thought，从而改善 theorem proving。 & 这篇 paper 构成本讲第 2 节的主轴，并且和课程早先关于 inference-time reasoning 的讨论直接相连。 \\
Chapter 10 & ImProver: Agent-Based Automated Proof Optimization\newline\url{https://arxiv.org/abs/2410.04753} & 在证明已经正确的前提下，能否用 agent 式流程优化 proof 的长度、可读性和模块化程度。 & 虽然 slides 主要讲 proving，而非 optimization，但它补上了“agent 不仅能找 proof，也能迭代改写 proof”的一环。 \\
Chapter 11 & An In-Context Learning Agent for Formal Theorem-Proving\newline\url{https://arxiv.org/abs/2310.04353} & 如果没有大量 environment-specific finetuning data，能否只靠强 LLM、proof environment 反馈和 search history 做 formal theorem proving。 & 这是数学 discovery 部分的主 reading，直接对应本讲对 theorem-proving agent 的讲解。 \\
Chapter 11 & Symbolic Regression with a Learned Concept Library\newline\url{https://arxiv.org/abs/2409.09359} & LLM 能否通过诱导和演化抽象 textual concepts，系统性改善 symbolic regression 的搜索效率与发现质量。 & 它支撑 lecture 下半场的 abstraction/discovery 主题，说明 agent 不只会证明，还会发明可复用的概念性搜索偏置。 \\
Chapter 12 & Privtrans: Automatically Partitioning Programs for Privilege Separation\newline\url{https://dawnsong.io/papers/privtrans.pdf} & How can a program be automatically partitioned so that privileged operations are confined to a smaller, more securable monitor while the rest runs with lower privilege? & This reading gives the lecture historical grounding for privilege separation: long before LLM agents, system security already relied on shrinking privilege boundaries. The lecture reuses the same principle for agent decomposition and tool sandboxing. \\
Chapter 12 & DataSentinel: A Game-Theoretic Detection of Prompt Injection Attacks\newline\url{https://arxiv.org/abs/2504.11358} & How can an LLM-based detector reliably identify prompt-injected inputs, including adaptive attacks that evolve to evade existing detectors? & This reading corresponds to the lecture's monitoring/detection layer. It is not positioned as a universal patch, but as one component in defense-in-depth for agent pipelines exposed to untrusted content. \\
Chapter 12 & AgentPoison: Red-teaming LLM Agents via Poisoning Memory or Knowledge Bases\newline\url{https://arxiv.org/abs/2407.12784} & What happens when an attacker poisons an agent's long-term memory or RAG knowledge base instead of injecting malicious instructions directly into the active prompt? & This reading grounds the lecture's warning that prompt injection is only one part of the attack surface. Memory poisoning and knowledge-base poisoning make the boundary problem persistent across sessions and tasks. \\
Chapter 12 & Progent: Programmable Privilege Control for LLM Agents\newline\url{https://arxiv.org/html/2504.11703v1} & How can an LLM agent be constrained so it only performs the tool calls necessary for the user task, while still preserving useful autonomy? & This is the lecture's central positive defense result. It turns least privilege from a principle into an enforceable runtime gate over tool calls, which is exactly where many agentic attacks become consequential. \\
\bottomrule
\end{longtable}
\normalsize
